\chapter{Ekuacionet diferenciale dhe dinamika numerike}
\begin{qellime}
Studenti zbaton Eulerin, Runge--Kutta dhe skemat me shumë hapa; analizon gabimin lokal, global dhe stabilitetin; simulon rënien e lirë, lavjerrësin dhe problemin e dy trupave; dallon probleme fillestare e kufitare.
\end{qellime}

\section{Problemi i vlerës fillestare}
Për
\begin{equation}
 \dot{\vect y}=\vect f(t,\vect y),\qquad \vect y(t_0)=\vect y_0,
\end{equation}
metoda e Eulerit është
\begin{equation}
 \vect y_{n+1}=\vect y_n+h\vect f(t_n,\vect y_n).
\end{equation}
Gabimi lokal është $\Oh(h^2)$ dhe gabimi global $\Oh(h)$. Euler-i është i dobishëm pedagogjikisht, por rrallë është zgjedhja përfundimtare.

Metoda e mesit dhe Heun janë të rendit të dytë. Metoda klasike RK4 përdor katër pjerrësi,
\begin{align}
 \vect k_1&=\vect f(t_n,\vect y_n),\\
 \vect k_2&=\vect f(t_n+h/2,\vect y_n+h\vect k_1/2),\\
 \vect k_3&=\vect f(t_n+h/2,\vect y_n+h\vect k_2/2),\\
 \vect k_4&=\vect f(t_n+h,\vect y_n+h\vect k_3),\\
 \vect y_{n+1}&=\vect y_n+\frac h6(\vect k_1+2\vect k_2+2\vect k_3+\vect k_4).
\end{align}
Çiftet e futura Runge--Kutta me dy rende japin vlerësim të gabimit dhe hap adaptiv.

\section{Stabiliteti}
Për ekuacionin provë $y'=\lambda y$, Euler-i jep $y_{n+1}=(1+h\lambda)y_n$. Stabiliteti kërkon $|1+h\lambda|<1$. Një metodë mund të ketë rend të lartë dhe përsëri të jetë e papërshtatshme për problem të ngurtë. Ngurtësia shfaqet kur shkallë kohore shumë të ndryshme imponojnë hap shumë të vogël për metoda eksplicite.

\section{Metodat me shumë hapa}
Adams--Bashforth përdor pjerrësi të mëparshme, për shembull
\begin{equation}
 y_{n+1}=y_n+\frac h2(3f_n-f_{n-1}).
\end{equation}
Këto metoda janë ekonomike pas nisjes, por kërkojnë histori dhe kujdes kur ndryshon hapi. Metodat implícite Adams--Moulton dhe BDF kanë zona stabiliteti më të favorshme për probleme të ngurta.

\section{Dinamikë Hamiltoniane}
Në mekanikë, ruajtja afatgjatë e strukturës mund të jetë më e rëndësishme se gabimi lokal minimal. Skema leapfrog/Verlet është simplektike:
\begin{align}
 \vect v_{n+1/2}&=\vect v_n+\frac h2\vect a(\vect r_n),\\
 \vect r_{n+1}&=\vect r_n+h\vect v_{n+1/2},\\
 \vect v_{n+1}&=\vect v_{n+1/2}+\frac h2\vect a(\vect r_{n+1}).
\end{align}
Ajo është e kthyeshme në kohë dhe mban energjinë të kufizuar rreth vlerës së saktë, në vend që të krijojë drift sistematik.

\begin{figure}[ht]
\centering\includegraphics[width=.74\textwidth]{10_energjia_lavjerresit.pdf}
\caption{Euler-i krijon drift energjie, ndërsa Verlet ruan më mirë strukturën afatgjatë.}
\end{figure}

\section{Shembuj fizikë}
\subsection{Rënia në ajër}
Për bosht pozitiv lart,
\begin{equation}
 \dot y=v,\qquad m\dot v=-mg-c_1v-c_2v|v|.
\end{equation}
Modeli lejon krahasimin e rezistencës lineare dhe kuadratike, shpejtësisë terminale dhe ndjeshmërisë ndaj parametrave.

\subsection{Lavjerrësi jolinear}
\begin{equation}
 \ddot\theta+\gamma\dot\theta+\frac gL\sin\theta=A\cos\Omega t.
\end{equation}
Pa forcim dhe fërkim, energjia ruhet. Me forcim mund të shfaqen rezonancë, shumëstabilitet dhe kaos; hapi duhet kontrolluar ndaj periodës më të shkurtër.

\subsection{Problemi i dy trupave}
Në njësi të shkallëzuara,
\begin{equation}
 \ddot{\vect r}=-\mu\frac{\vect r}{r^3}.
\end{equation}
Kontrollet janë energjia, momenti këndor dhe ligji i tretë i Keplerit. Një orbitë që duket e mbyllur për pak perioda mund të ketë drift të ngadalshëm numerik.

\begin{figure}[ht]
\centering\includegraphics[width=.60\textwidth]{11_orbita_dy_trupave.pdf}
\caption{Orbitë e integruar me skemë simplektike.}
\end{figure}

\section{Problemet e vlerave kufitare}
Në një problem kufitar, kushtet jepen në skaje të ndryshme. Metoda e qitjes e kthen problemin në një kërkim të kushtit fillestar që plotëson skajin tjetër. Diferencat e fundme e kthejnë ekuacionin në sistem algjebrik. Qitja mund të jetë e ndjeshme; formulimi matricor është shpesh më robust.

\begin{lidhje}
Laboratori krahason Euler, RK4 dhe Verlet për rënien, lavjerrësin dhe orbitën. Për çdo metodë vizatohen gabimi, energjia dhe kostoja kundrejt hapit.
\end{lidhje}

\section*{Pyetje kontrolli}
\begin{enumerate}
\item Dalloni rendin nga stabiliteti.
\item Pse një metodë simplektike mund të jetë më e mirë për orbita afatgjata?
\item Si verifikohet numerikisht një integrues orbital?
\end{enumerate}
\section{Integruesit kohorë në Python}
\subsection{Euler dhe Runge--Kutta i rendit të katërt}
\begin{lstlisting}[language=Python,caption={Dy integrues të përgjithshëm për sisteme ODE.}]
def euler(f, t, y0):
    y = np.empty((len(t), len(np.atleast_1d(y0))), float)
    y[0] = y0
    for n, h in enumerate(np.diff(t)):
        y[n+1] = y[n] + h*f(t[n], y[n])
    return y

def rk4(f, t, y0):
    y = np.empty((len(t), len(np.atleast_1d(y0))), float)
    y[0] = y0
    for n, h in enumerate(np.diff(t)):
        k1 = f(t[n], y[n])
        k2 = f(t[n]+h/2, y[n]+h*k1/2)
        k3 = f(t[n]+h/2, y[n]+h*k2/2)
        k4 = f(t[n]+h, y[n]+h*k3)
        y[n+1] = y[n] + h*(k1+2*k2+2*k3+k4)/6
    return y
\end{lstlisting}

\subsection{Modeli i lavjerrësit}
\begin{lstlisting}[language=Python]
def lavjerresi(t, y, g=9.81, L=1.0, gamma=0.0):
    theta, omega = y
    return np.array([omega, -g/L*np.sin(theta) - gamma*omega])

def energjia(y, g=9.81, L=1.0):
    theta, omega = y.T
    return 0.5*(L*omega)**2 + g*L*(1-np.cos(theta))
\end{lstlisting}
Për rastin pa fërkim, drift-i i energjisë mat cilësinë afatgjatë. Për rastin me fërkim, energjia duhet të bjerë; ruajtja nuk është më kontrolli i duhur.

\subsection{Velocity Verlet për gravitetin}
\begin{lstlisting}[language=Python,caption={Skema simplektike për orbitën.}]
def verlet_orbite(r0, v0, dt, n, mu=1.0):
    r = np.empty((n, 2)); v = np.empty((n, 2))
    r[0], v[0] = r0, v0
    def a(q):
        return -mu*q/np.linalg.norm(q)**3
    for k in range(n-1):
        v_half = v[k] + 0.5*dt*a(r[k])
        r[k+1] = r[k] + dt*v_half
        v[k+1] = v_half + 0.5*dt*a(r[k+1])
    return r, v
\end{lstlisting}

\subsection{Protokolli i verifikimit}
\begin{enumerate}
\item Euler-i testohet te $y'=\lambda y$, ku zgjidhja njihet.
\item RK4 testohet me përgjysmim hapi dhe pritet faktor gabimi afërsisht $2^4$.
\item Për orbitën kontrollohen energjia dhe momenti këndor.
\item Për lavjerrësin e këndit të vogël krahasohet perioda me $2\pi\sqrt{L/g}$.
\item Koha dhe memoria raportohen bashkë me gabimin.
\end{enumerate}

\begin{kujdes}
Në terminologjinë numerike shqipe përdoren \emph{metodë e shtjelluar} dhe \emph{metodë e pashtjelluar} për explicit dhe implicit. Në tekst jepen edhe termat ndërkombëtarë kur ndihmojnë studentin të lexojë dokumentacionin e bibliotekave.
\end{kujdes}

Për terminologjinë \emph{problem i vlerës fillestare}, \emph{gabim lokal}, \emph{gabim global}, \emph{metodë e shtjelluar} dhe \emph{qëndrueshmëri} janë përdorur gjerësisht burimet shqip \cite{dishmema,hanelli,hoxha,datja,dishmemaArticle}.

\subsection{Gabimi lokal dhe global i Eulerit}
Për problemin $\vect y'=\vect f(t,\vect y)$, Taylor-i i zgjidhjes së saktë është
\begin{equation}
 \vect y(t_n+h)=\vect y(t_n)+h\vect f(t_n,\vect y(t_n))
 +\frac{h^2}{2}\vect y''(\xi_n).
\end{equation}
Euler-i heq termin e fundit; defekti lokal është $\Oh(h^2)$. Nëse $\vect f$ është Lipschitz në $\vect y$ me konstante $L$, gabimet plotësojnë afërsisht
\begin{equation}
 \|\vect e_{n+1}\|\le(1+hL)\|\vect e_n\|+Ch^2.
\end{equation}
Pas $N=T/h$ hapash dhe përdorimit të pabarazisë diskrete të Grönwall-it merret $\|\vect e_N\|\le C_T h$: rendi global është një më i ulët se rendi i defektit lokal.

\subsection{Qëndrueshmëria absolute}
Ekuacioni provë $y'=\lambda y$ ndan metodën nga problemi. Euler-i jep
\begin{equation}
 y_{n+1}=(1+h\lambda)y_n.
\end{equation}
Për një zgjidhje fizike që shuhet, $\Re\lambda<0$, kërkohet $|1+h\lambda|<1$. Bashkësia $\{z:|1+z|<1\}$ është rajoni i qëndrueshmërisë absolute. Për $\lambda<0$ real, kushti bëhet $0<h<2/|\lambda|$. Një hap mund të jetë mjaft i vogël për saktësi dhe përsëri i papërshtatshëm për qëndrueshmëri, ose anasjelltas.

Problemet e shtyllura përmbajnë shkallë kohore shumë të ndryshme. Metodat eksplicite detyrohen të ndjekin shkallën më të shpejtë edhe kur interesi fizik është te evolucioni i ngadaltë. Euler-i i prapëm ka faktor $1/(1-h\lambda)$ dhe është i qëndrueshëm për të gjithë gjysmëplanin $\Re(h\lambda)<0$, por kërkon zgjidhje algjebrike në çdo hap.

\subsection{Runge--Kutta i rendit të katërt}
Metodat Runge--Kutta përputhin zhvillimin Taylor pa llogaritur derivate të larta të $\vect f$. Skema klasike është
\begin{align}
 \vect k_1&=\vect f(t_n,\vect y_n),\\
 \vect k_2&=\vect f(t_n+h/2,\vect y_n+h\vect k_1/2),\\
 \vect k_3&=\vect f(t_n+h/2,\vect y_n+h\vect k_2/2),\\
 \vect k_4&=\vect f(t_n+h,\vect y_n+h\vect k_3),\\
 \vect y_{n+1}&=\vect y_n+\frac{h}{6}(\vect k_1+2\vect k_2+2\vect k_3+\vect k_4).
\end{align}
Koeficientët plotësojnë kushtet e rendit që barazojnë zhvillimin e skemës me Taylor-in deri në $h^4$. Gabimi lokal është $\Oh(h^5)$ dhe globali $\Oh(h^4)$. Katër vlerësime të funksionit kushtojnë më shumë se Euler-i, por për një tolerancë të dhënë zakonisht lejojnë hapa shumë më të mëdhenj.

\subsection{Metodat me shumë hapa}
Integrimi i identitetit
\begin{equation}
 \vect y(t_{n+1})=\vect y(t_n)+\int_{t_n}^{t_{n+1}}\vect f(t,\vect y(t))\dd t
\end{equation}
dhe interpolimi i $\vect f$ në pika të mëparshme jep Adams--Bashforth. Për dy hapa,
\begin{equation}
 \vect y_{n+1}=\vect y_n+h\left(\frac32\vect f_n-\frac12\vect f_{n-1}\right).
\end{equation}
Metoda kursen vlerësime të funksionit, por nuk nis vetë: nevojitet një metodë njëhapëshe për gjendjet fillestare. Stabiliteti dhe rrënjët e polinomit karakteristik duhet të kontrollohen; një formulë me rend formal të lartë mund të jetë zero-e-paqëndrueshme.

\subsection{Leapfrog, Verlet dhe dinamika Hamiltoniane}
Për $\ddot q=a(q)$, velocity Verlet shkruhet
\begin{align}
 v_{n+1/2}&=v_n+\frac h2a(q_n),\\
 q_{n+1}&=q_n+hv_{n+1/2},\\
 v_{n+1}&=v_{n+1/2}+\frac h2a(q_{n+1}).
\end{align}
Skema është e kthyeshme në kohë dhe simplektike. Ajo nuk e ruan energjinë pikë për pikë, por ruan një Hamiltonian të modifikuar pranë atij të vërtetë; prandaj gabimi i energjisë lëkundet pa drift sistematik për kohë të gjata. Kjo është arsyeja strukturore pse Verlet-i mund të jetë më i mirë se një metodë me rend lokal më të lartë për orbita afatgjata.

\subsection{Problemet kufitare}
Për $y''=f(x,y,y')$ me $y(a)=\alpha$, $y(b)=\beta$, metoda e qitjes e zëvendëson pjerrësinë e panjohur $s=y'(a)$ me parametër. Integrohet problemi fillestar dhe zgjidhet ekuacioni skalar
\begin{equation}
 \Phi(s)=y(b;s)-\beta=0.
\end{equation}
Qitja është intuitive, por mund të jetë e keqkushtëzuar. Diferencat e fundme ndërtojnë një sistem algjebrik për të gjitha nyjet dhe janë më të përshtatshme kur zgjidhja rritet ose shuhet shumë shpejt.

\subsection{Përmbledhje dhe ushtrime}
Zgjedhja e integruesit varet jo vetëm nga rendi, por nga stabiliteti, struktura e sistemit dhe horizonti kohor.
\begin{enumerate}
\item Nxirrni rajonin e qëndrueshmërisë së Eulerit të përmirësuar.
\item Verifikoni rendet e Eulerit dhe RK4 me metodën e përgjysmimit të hapit.
\item Krahasoni RK4 dhe Verlet për orbitën e Keplerit duke matur energjinë dhe momentin këndor.
\item Zgjidhni një problem kufitar me qitje dhe me diferenca të fundme; diskutoni kushtëzimin.
\end{enumerate}

\subsection{Lavjerrësi jolinear dhe perioda e varur nga amplituda}
Ekuacioni pa fërkim është
\begin{equation}
 \ddot\theta+\frac{g}{L}\sin\theta=0.
\end{equation}
Linearizimi $\sin\theta\simeq\theta$ jep periodën $T_0=2\pi\sqrt{L/g}$, por modeli linear humbet varësinë nga amplituda. Nga ruajtja e energjisë,
\begin{equation}
 \frac12L^2\dot\theta^2+gL(1-\cos\theta)
 =gL(1-\cos\theta_0),
\end{equation}
merret perioda
\begin{equation}
 T=4\sqrt{\frac{L}{g}}
 \int_0^{\pi/2}\frac{\dd\phi}
 {\sqrt{1-k^2\sin^2\phi}},
 \qquad k=\sin\frac{\theta_0}{2}.
\end{equation}
Ky integral eliptik mund të llogaritet me kuadraturë dhe të krahasohet me periodën e nxjerrë nga integrimi kohor. Për amplituda të vogla,
\begin{equation}
 \frac{T}{T_0}=1+\frac{\theta_0^2}{16}
 +\frac{11\theta_0^4}{3072}+\cdots.
\end{equation}
Krahasimi i tri rrugëve --- seria, kuadratura dhe ODE-ja --- është verifikim i fuqishëm ndërmetodik.

Perioda numerike nga trajektorja matet me kalime në zero ose me kulme. Interpolimi ndërmjet dy hapave zvogëlon gabimin e kuantizimit kohor. Për integrim afatgjatë, monitorohet energjia. Një metodë jo simplektike mund të japë periodë të mirë për disa cikle dhe pastaj të krijojë drift që ndryshon artificialisht amplitudën.

\subsubsection{Problemi i dy trupave dhe reduktimi në një trup}
Për masa $m_1,m_2$ me pozicione $\vect r_1,\vect r_2$, forcat janë të barabarta e të kundërta. Qendra e masës
\begin{equation}
 \vect R=\frac{m_1\vect r_1+m_2\vect r_2}{m_1+m_2}
\end{equation}
lëviz në mënyrë uniforme, ndërsa koordinata relative $\vect r=\vect r_1-\vect r_2$ plotëson
\begin{equation}
 \mu\ddot{\vect r}=-\frac{Gm_1m_2}{r^3}\vect r,
 \qquad \mu=\frac{m_1m_2}{m_1+m_2}.
\end{equation}
Pas pjesëtimit me $\mu$ merret $\ddot{\vect r}=-G(m_1+m_2)\vect r/r^3$. Invariantet janë energjia
\begin{equation}
 E=\frac12\mu|\dot{\vect r}|^2-\frac{Gm_1m_2}{r}
\end{equation}
dhe momenti këndor $\vect L=\mu\vect r\times\dot{\vect r}$. Për orbitë të lidhur me gjysmëbosht të madh $a$,
\begin{equation}
 T^2=\frac{4\pi^2}{G(m_1+m_2)}a^3.
\end{equation}
Një laborator i plotë nxjerr periodën nga simulimi për disa $a$, përshtat pjerrësinë në shkallë logaritmike dhe kontrollon konstanten fizike, jo vetëm eksponentin $3/2$.

\subsubsection{Gabimi në ngjarje dhe hapa adaptivë}
Shumë probleme kërkojnë kohën e një ngjarjeje: goditja e tokës, kalimi në periapsis ose arritja e një pragu. Nëse ngjarja ndodh ndërmjet $t_n$ dhe $t_{n+1}$, kthimi i thjeshtë i $t_{n+1}$ ka gabim të rendit $h$. Duhet interpoluar funksioni i ngjarjes $g(t,y)$ dhe zgjidhur $g=0$ brenda hapit. Bibliotekat ODE e bëjnë këtë me interpolues të dendur, por përdoruesi duhet të përcaktojë drejtimin e kalimit dhe nëse integrimi ndalet.

Kontrolli adaptiv përdor dy vlerësime me rende të ndryshme ose step-doubling. Nëse gabimi lokal i vlerësuar është $e$, hapi i ri zgjidhet afërsisht
\begin{equation}
 h_{\mathrm{ri}}=s h\left(\frac{\mathrm{tol}}{e}\right)^{1/(p+1)},
\end{equation}
me faktor sigurie $s<1$. Toleranca absolute mbron komponentët pranë zeros; toleranca relative shkallëzon komponentët e mëdhenj. Në sisteme me njësi të ndryshme nevojiten toleranca komponentore ose nondimensionalizim.

